% ============================================================================
%  Section V — Methodology  (IOP MST)
%  Passive voice / third person throughout.
%  Requires: booktabs, amsmath, tikz + \usetikzlibrary{arrows.meta,positioning}.
% ============================================================================

\section{Methodology}
\label{sec:method}

The processing chain evaluated in this work is deliberately lightweight, so that
it is representative of what an inexpensive sensor node can execute: fixed-length
windowing, a compact set of time-domain features, and a conventional classifier.
The chain is illustrated in Fig.~\ref{fig:pipeline}. The emphasis of this section,
and the principal methodological contribution of the paper, is the set of
\emph{evaluation protocols} used to obtain deployment-realistic performance
estimates.

% ---------------------------------------------------------------------------
\begin{figure}[t]
\centering
\definecolor{pipeBlue}{HTML}{0072B2}
\definecolor{pipeSky}{HTML}{56B4E9}
\definecolor{pipeGreen}{HTML}{009E73}
\definecolor{pipeAmber}{HTML}{E69F00}
\definecolor{pipeRed}{HTML}{D55E00}
\begin{tikzpicture}[
  font=\footnotesize,
  >={Latex[length=1.7mm,width=1.3mm]},
  stage/.style={draw=#1!80!black,fill=#1!6,line width=0.6pt,rounded corners=2.5pt,
                align=center,inner sep=3pt,text width=2.50cm,minimum height=1.42cm},
  stage/.default=black,
  flow/.style={->,line width=0.7pt,draw=black!45},
  chip/.style={draw=#1!80!black,fill=#1!7,line width=0.5pt,rounded corners=2.5pt,
               align=center,inner sep=3pt,text width=3.00cm,minimum height=0.95cm},
  chip/.default=black
]

% ---- processing chain ---------------------------------------------------
\node[stage=pipeBlue] (sens)
  {\textcolor{pipeBlue}{\bfseries Smartphone}\\[-1pt]
   \textcolor{pipeBlue}{\bfseries MEMS node}\\[1.5pt]
   {\scriptsize\color{black!58}$100$~Hz, $6$ channels}};
\node[stage=pipeSky,right=0.26cm of sens] (win)
  {\textcolor{pipeSky!75!black}{\bfseries Windowing}\\[1.5pt]
   {\scriptsize\color{black!58}$2$~s windows}\\[-1pt]
   {\scriptsize\color{black!58}$50\%$ overlap}};
\node[stage=pipeSky,right=0.26cm of win] (feat)
  {\textcolor{pipeSky!75!black}{\bfseries Time-domain}\\[-1pt]
   \textcolor{pipeSky!75!black}{\bfseries features}\\[1.5pt]
   {\scriptsize\color{black!58}$14$ per channel}};
\node[stage=pipeGreen,right=0.26cm of feat] (clf)
  {\textcolor{pipeGreen!80!black}{\bfseries Baseline}\\[-1pt]
   \textcolor{pipeGreen!80!black}{\bfseries classifiers}\\[1.5pt]
   {\scriptsize\color{black!58}RF / KNN / SVM}};
\node[stage=pipeAmber,right=0.26cm of clf] (out)
  {\textcolor{pipeAmber!70!black}{\bfseries Health state}\\[1.5pt]
   {\scriptsize\color{black!58}one of six}\\[-1pt]
   {\scriptsize\color{black!58}H, B1--B3, V, R}};

\draw[flow] (sens)--(win);
\draw[flow] (win)--(feat);
\draw[flow] (feat)--(clf);
\draw[flow] (clf)--(out);

% ---- evaluation-protocol panel -----------------------------------------
\node[draw=black!35,fill=black!3,rounded corners=3pt,line width=0.5pt,
      minimum width=14.60cm,minimum height=1.98cm,below=0.70cm of feat] (panel) {};

\node[anchor=north west,font=\scriptsize,color=black!70]
  at ([xshift=0.32cm,yshift=-0.16cm]panel.north west)
  {\textbf{Evaluation protocol}\,---\,how the windows are partitioned into
   training and test sets};

\node[chip=pipeRed,anchor=south west]
  at ([xshift=0.34cm,yshift=0.20cm]panel.south west) (c1)
  {\textcolor{pipeRed}{\bfseries RANDOM}\\[1pt]
   {\scriptsize\color{black!58}random window split}\\[-1pt]
   {\scriptsize\itshape\color{pipeRed}leaky reference}};
\node[chip=pipeGreen,right=0.22cm of c1] (c2)
  {\textcolor{pipeGreen!80!black}{\bfseries GROUP}\\[1pt]
   {\scriptsize\color{black!58}recordings held out}\\[-1pt]
   {\scriptsize\itshape\color{pipeGreen!80!black}leakage-free}};
\node[chip=pipeGreen,right=0.22cm of c2] (c3)
  {\textcolor{pipeGreen!80!black}{\bfseries CROSS-SPEED}\\[1pt]
   {\scriptsize\color{black!58}unseen supply frequency}\\[-1pt]
   {\scriptsize\itshape\color{pipeGreen!80!black}$+$ speed shift}};
\node[chip=pipeGreen,right=0.22cm of c3] (c4)
  {\textcolor{pipeGreen!80!black}{\bfseries CROSS-LOAD}\\[1pt]
   {\scriptsize\color{black!58}unseen load condition}\\[-1pt]
   {\scriptsize\itshape\color{pipeGreen!80!black}$+$ load shift}};

\draw[dashed,line width=0.6pt,draw=black!45,->] (feat.south) -- (panel.north);

\end{tikzpicture}
\caption{Processing chain and evaluation protocols. Identical features are
assessed under four data-partitioning protocols: RANDOM is retained only as a
leaky reference, whereas GROUP, CROSS-SPEED and CROSS-LOAD are leakage-free at
the recording level, the latter two additionally imposing an operating-condition
shift. The protocol, not the model, is the primary object of study.}
\label{fig:pipeline}
\end{figure}

\subsection{Segmentation and Windowing}
Let the dataset comprise $R=36$ continuous recordings, each associated with a
health label and an operating condition (supply frequency and load). Every
recording was divided into fixed-length windows of $2$~s with $50\%$ overlap,
matching the segmentation of the source data article. At $100$~Hz this yields
windows of $200$~samples and a total of $32{,}328$ windows, distributed equally
across the six classes ($5{,}388$ per class). The window length and overlap were
treated as configurable parameters and were varied in the trade-off study of
Section~\ref{sec:results}.

\subsection{Feature Extraction}
Each window was reduced to a set of $14$ time-domain features per channel,
summarised in Table~\ref{tab:features}. These comprise standard statistical and
shape descriptors widely used in vibration-based diagnosis, and require no
spectral transform, which keeps the extraction cost compatible with a
microcontroller. For a window of $C$ channels, the per-channel feature vectors
were concatenated into a single descriptor of length $14C$ (for example, $84$ for
the full six-channel set). The six-feature subset marked in
Table~\ref{tab:features}, corresponding to the descriptors used by the source
data article's baseline, was retained for a feature-set ablation. For the
distance- and margin-based classifiers, features were standardised to zero mean
and unit variance using statistics estimated on the training partition only.

\subsection{Classifiers}
Three baseline classifiers were considered, chosen for their prevalence in the
fault-diagnosis literature and their modest computational footprint: a Random
Forest (RF) with $200$ trees; a $k$-nearest-neighbour classifier (KNN, $k=5$);
and a support vector machine (SVM) with a radial-basis-function kernel. To keep
the SVM tractable on the full window set, the kernel was approximated by a
Nystr\"om feature map ($400$ components) followed by a linear SVM. For the
edge-footprint analysis of Section~\ref{sec:results}, two additional
lightweight models---a depth-limited decision tree and a multinomial logistic
regression---were included. All models map a feature descriptor to one of the six
health states.

\subsection{Evaluation Protocols}
\label{subsec:protocols}
The central methodological question is how the labelled windows are partitioned
into training and test sets. Four protocols were defined and applied to identical
features so that differences in reported accuracy are attributable solely to the
partitioning strategy.

\subsubsection{RANDOM}
A stratified random split over all windows, as commonly used in the literature
and by the source baseline. This protocol is retained as a reference precisely
because it is susceptible to leakage: adjacent windows share $50\%$ of their
samples, and all windows of a recording share the same operating point and sensor
coupling, so a model may recognise the originating recording rather than the
fault.

\subsubsection{GROUP}
A leakage-free protocol in which entire recordings are held out. Stratified
group $k$-fold cross-validation ($k=6$) was used, with the recording identifier as
the grouping variable, so that no window from a test recording appears in
training. With six recordings per class, each fold holds out one recording per
class.

\subsubsection{RANDOM without overlap and BLOCKED}
Two further partitionings are defined to identify which property of the RANDOM
protocol is responsible for its optimism. The first repeats the stratified random
split on windows extracted with zero overlap, so that no sample can appear in both
partitions. The second holds out a contiguous temporal block from within every
recording, taking successive sixths of each recording in turn, so that training
and test windows are additionally separated in time while every recording still
contributes to training. Comparing these against RANDOM and GROUP attributes the
gap to shared samples, to temporal proximity, or to recording identity.

\subsubsection{CROSS-SPEED}
A leave-one-supply-frequency-out protocol: the model is trained on two supply
frequencies and tested on the third, and results are averaged over the three
folds. This quantifies generalisation to an unseen operating speed.

\subsubsection{CROSS-LOAD}
A leave-one-load-out protocol: the model is trained on one load condition
(loaded or unloaded) and tested on the other. This quantifies generalisation to
an unseen load.

The GROUP, CROSS-SPEED, and CROSS-LOAD partitions are all group-disjoint at the
recording level and therefore free of the window-overlap leakage that affects the
RANDOM protocol; the latter two additionally introduce a deliberate operating
-condition shift between training and test.

\subsection{Reporting performance and its dispersion}
\label{subsec:stats}
Classification performance is reported as accuracy and macro-averaged
$F_1$-score, the latter being sensitive to the per-class failures that a balanced
accuracy can hide. Two questions then arise: over what unit is performance
replicated, and how is that dispersion to be described.

The unit adopted here is the \emph{held-out recording}. A score is computed
separately for each recording while it is held out, so the reported mean is an
average over $36$ recordings and its dispersion is the recording-to-recording
variation, which is the variation a new installation would actually encounter.
Intervals are $95\%$ intervals for that mean, clustered on recording. Ten random
seeds are used, each reshuffling the recording-level fold composition, but the
seeds serve only to stabilise the per-recording means; they are not treated as
independent observations, because re-partitioning the same $36$ recordings
supplies no new information about the machine and an interval computed over seeds
understates the true dispersion by a large factor. In the present data the
recording-clustered interval is roughly three times the seed-level interval, which
is the size of the error that the more convenient choice would introduce.

For the same reason no claim of a Type~A evaluation of measurement uncertainty in
the sense of the GUM~\cite{jcgm2008gum} is made for these figures, although that
framework is the natural destination for such work~\cite{angrisani2025uncertainty,
dorozhovets2020uncertainty}. A GUM-conformant budget would require the dominant
influence quantities---machine, fault specimen, mounting, session---to be varied
and their contributions evaluated~\cite{staszewski2025mems}; none of them varies
within this dataset. What is reported is therefore a reproducibility interval over
the available recordings, and the influence quantities that remain unquantified
are set out qualitatively in Section~\ref{subsec:limits}.

Configurations are compared on their per-recording differences with a two-sided
paired permutation test over the $2^{n}$ sign changes, approximated by $20\,000$
random sign assignments. The test is exact in construction, assumes neither
normality nor equal variances, and---unlike a signed-rank test over five
seeds---does not exhaust its resolution before reaching conventional significance
levels. Confusion matrices are aggregated from the pooled out-of-fold predictions
of the GROUP protocol.

Finally, where a configuration is selected rather than fixed in advance, selection
and estimation are separated. Choosing an operating point by inspecting
cross-validation scores and then quoting those same scores is optimistically
biased~\cite{qi2019cv}; Section~\ref{subsec:nested} therefore repeats the
selection inside a nested recording-wise loop, so that the recordings used to
estimate performance are never seen by the selector.

\subsection{Implementation Details}
\label{subsec:impl}
All estimators were taken from a single, standard machine-learning library, with
the settings collected in Table~\ref{tab:impl}; any parameter not listed was left
at its library default, and every stochastic component was seeded. Three further
choices affect reproducibility and are therefore stated explicitly. First,
decimation to a lower sampling rate was performed by polyphase resampling with
the associated anti-alias FIR filter, so that the reduced-rate conditions of
Section~\ref{subsec:tradeoff} combine a change of rate with the low-pass that
necessarily accompanies it. Second, the mutual information of
Section~\ref{sec:charact} was obtained with a $k$-nearest-neighbour estimator
($k=3$) applied to the continuous features and reported in nats; the estimator is
non-negative by construction, so values near zero indicate no detectable
dependence rather than independence. Third, the power spectral densities were
computed by Welch's method with a Hann window, $50\%$ segment overlap and
density scaling, using $2048$-sample segments for the per-class spectra of
Fig.~\ref{fig:psd} and $1024$-sample segments for the signal-floor estimate.

% ---------------------------------------------------------------------------
\begin{table}[t]
\centering
\caption{Estimator and feature-extraction settings. Parameters not listed were
left at their library defaults; $d$ denotes the feature-vector dimension
($14C$ for $C$ channels).}
\label{tab:impl}
\renewcommand{\arraystretch}{1.25}
\small
\begin{tabular}{@{}p{3.6cm}p{10.9cm}@{}}
\toprule
\textbf{Component} & \textbf{Setting} \\
\midrule
Random Forest (baseline)   & $200$ trees, unlimited depth, Gini split,
                             $\sqrt{d}$ features per split \\
Random Forest (edge)       & $50$ trees, maximum depth $8$ \\
Decision tree (edge)       & maximum depth $8$ \\
KNN                        & $k=5$, Euclidean metric, uniform weights \\
SVM                        & RBF kernel approximated by a Nystr\"om map
                             ($400$ components, $\gamma=1/d$), then a linear
                             SVM ($C=10$, squared hinge, primal solver,
                             $\leq5000$ iterations) \\
Logistic regression        & multinomial, $L_2$ penalty, $C=1$, L-BFGS,
                             $\leq2000$ iterations \\
Standardisation            & zero mean, unit variance; fitted on the training
                             partition only (KNN, SVM, logistic regression) \\
Amplitude entropy          & $32$-bin normalised histogram per window \\
Bit-depth reduction        & per-recording, per-channel min--max range
                             uniformly quantised to $2^{b}-1$ levels \\
Packet loss                & independent Bernoulli per sample, concealed by
                             sample-and-hold \\
Orientation change         & Rodrigues rotation of both triads by
                             $30^{\circ}$ about the axis $(1,1,0)/\sqrt{2}$ \\
\bottomrule
\end{tabular}
\end{table}

% ---------------------------------------------------------------------------
\begin{table}[t]
\centering
\caption{Per-channel time-domain features. For a window
$x=\{x_i\}_{i=1}^{W}$ with mean $\mu$ and standard deviation $\sigma$.
Features marked $\dagger$ form the six-feature ablation subset.}
\label{tab:features}
\renewcommand{\arraystretch}{1.45}
\begin{tabular}{@{}ll@{}}
\toprule
\textbf{Feature} & \textbf{Definition} \\
\midrule
Mean & $\mu=\frac{1}{W}\sum_i x_i$ \\
Std.\ dev.$^\dagger$ & $\sigma=\sqrt{\frac{1}{W}\sum_i (x_i-\mu)^2}$ \\
Variance & $\sigma^2$ \\
RMS$^\dagger$ & $x_{\mathrm{rms}}=\sqrt{\frac{1}{W}\sum_i x_i^2}$ \\
Peak-to-peak$^\dagger$ & $\max_i x_i-\min_i x_i$ \\
Max.\ absolute & $x_{\max}=\max_i |x_i|$ \\
Skewness$^\dagger$ & $\frac{1}{W}\sum_i (x_i-\mu)^3/\sigma^3$ \\
Kurtosis$^\dagger$ & $\frac{1}{W}\sum_i (x_i-\mu)^4/\sigma^4-3$ \\
Crest factor$^\dagger$ & $x_{\max}/x_{\mathrm{rms}}$ \\
Shape factor & $x_{\mathrm{rms}}/\big(\frac{1}{W}\sum_i |x_i|\big)$ \\
Impulse factor & $x_{\max}/\big(\frac{1}{W}\sum_i |x_i|\big)$ \\
Clearance factor & $x_{\max}/\big(\frac{1}{W}\sum_i \sqrt{|x_i|}\big)^2$ \\
Zero-crossing rate & rate of sign changes of $x_i-\mu$ \\
Amplitude entropy & $-\sum_k p_k\log_2 p_k$ (32-bin histogram) \\
\bottomrule
\end{tabular}
\end{table}
